Se busca una ecuación autónoma \(y'=f(y)\), donde \(f\) sea un polinomio real del menor grado posible. La recta de fases debe tener exactamente los equilibrios \(0\) y \(1\), con flechas hacia valores menores de \(y\) para \(y<0\) y hacia valores mayores para \(0<y<1\) y para \(y>1\). Además, en \(y=2\) la pendiente debe ser \(2\).

Construya la ecuación y justifique por qué no basta un polinomio de grado menor. Sin resolverla explícitamente, describa si las trayectorias cercanas a cada equilibrio se acercan o se alejan por cada lado. Para \(0<y(0)<1\), determine el comportamiento cuando \(t\to\infty\); puede usar que una solución contenida en un intervalo compacto donde el campo es suave se prolonga para todo tiempo positivo.
